% New structural comparison, separate from any blanket lower bound.
\subsection{Removing an opposite corner with projective coinvariants}

\begin{lemma}[Projective trivial modules and the minimal column]
\label{C:projective-coinvariants-minimal-column}
Let $N$ be a finite monoid with minimal Rees ideal $I\times G\times J$.
The trivial right $\mathbb QN$-module is projective if and only if
$|J|=1$. In this case a minimal idempotent $e$ satisfies
$eN=eNe=G$, and
\[
 q_N:N\longrightarrow G,\qquad n\longmapsto en
\]
is a unital monoid homomorphism, with $e$ the identity of $G$.
Over a field of characteristic $p$, the same trivial right module
is projective if and only if $|J|=1$ and $p\nmid|G|$.
\end{lemma}
\begin{proof}
A splitting of the right augmentation is an element $a\in kN$
with augmentation one and $an=a$ for every $n\in N$.
For each minimal idempotent $f$ it follows that $a=af\in kNf$.
Distinct minimal columns give disjoint basis supports for these
left ideals. Their intersection is zero, contradicting the
augmentation of $a$. Thus there is only one minimal column.

For such an ideal, $eN$ is the group $eNe$.
Consequently $(en)e=en$ and
$(en)(em)=enm$, proving the homomorphism assertion.
If $|G|$ is invertible in $k$, the element
$|G|^{-1}\sum_{g\in G}g$ is right invariant under $N$, since
right multiplication by $n$ acts on this sum as right multiplication
by $en\in G$. Its augmentation is one, giving the splitting.

Conversely, any invariant element belongs to $kK$, since $a=ae$.
Its coefficients on each fixed-row $G$-orbit are constant under
right multiplication by $G$. Its augmentation is therefore a
sum of multiples of $|G|$. In characteristic $p$ dividing $|G|$
this cannot be one.
\end{proof}

\begin{lemma}[The equivariant action-groupoid comparison]
\label{C:opposite-corner-action-groupoid}
Let $N$ satisfy the preceding lemma, and let $S$ be a left
$N$-set with a commuting right action by a monoid $M$.
Let $S\mathbin{/\mkern-6mu/}N$ denote its left action category.
There is an actual right-$M$-equivariant functor
\[
 p:S\mathbin{/\mkern-6mu/}N
       \longrightarrow eS\mathbin{/\mkern-6mu/}G,
 \qquad s\longmapsto es,\quad n\longmapsto en,
\]
which induces a weak equivalence of classifying spaces.
Moreover $N\backslash S$ is the finite right $M$-set
$T=G\backslash eS$.
\end{lemma}
\begin{proof}
The action category has an arrow labelled $n$ from $s$ to $ns$.
The displayed functor preserves its target because
$(en)(es)=ens$, and preserves compositions because $q_N$ is a
homomorphism. It commutes strictly with the right $M$-action.

For $t\in eS$, an object of $(t\downarrow p)$ is a pair $(s,g)$
with $g\in G$ and $es=gt$. The pair $(t,e)$ is one such object.
There is a morphism from $(s,g)$ to $(t,e)$ labelled $g^{-1}\in N$:
indeed $g^{-1}s=g^{-1}es=t$ and $e g^{-1}g=e$.
For an arrow labelled $n$ from $(s,g)$ to $(s',g')$, the comma
equation is $g'=(en)g$. Hence
\[
 (g')^{-1}n=(g')^{-1}(en)=g^{-1}.
\]
These arrows give a natural transformation from the identity of
the comma category to its constant functor at $(t,e)$.
The component at $(t,e)$ need not be the identity of $N$;
this is immaterial to the contraction of its nerve.
Thus every comma category is contractible, and Quillen's theorem A
proves the weak equivalence.

Every $s$ is related to $es$ under the $N$-orbit equivalence.
For $n\in N$, the states $ens=(en)(es)$ and $es$ lie in the same
$G$-orbit. This gives inverse identifications of the two orbit
sets, compatible with their right $M$-actions.
\end{proof}


\begin{theorem}[Finite-target and good-prime corner removal]
\label{C:finite-target-opposite-corner}
Let $C$ be the triangular two-object category with endomorphism
monoids $M,N$, cross-biset $S$ with left $N$-action and right
$M$-action, and no reverse cross arrows.
Assume the trivial right $\mathbb QN$-module is projective.
Put $T=N\backslash S$ and let
\[
 Q:C\longrightarrow C(M,T)
\]
be the actual quotient functor, equal to the identity on $M$,
sending $N$ to the identity at the other object and $S$ to $T$.
Then:
\begin{enumerate}
\item $BQ$ induces an isomorphism of cohomology rings over
$\mathbb Q$, and over $\mathbb F_p$ whenever $p\nmid|G|$.
\item For every connected finite CW complex $X$, precomposition is
a weak equivalence of derived unbased mapping spaces
\[
 \operatorname{Map}(BC(M,T),X)
       \longrightarrow \operatorname{Map}(BC,X).
\]
\item If $BC$ has finite CW homotopy type, there is a homotopy
left inverse $h:BC(M,T)\to BC$ to $BQ$.
\end{enumerate}
The size of the act category is $|M|+|T|+1\leq|M|+|S|+1$.
No assertion that it has finite CW type or the same integral
homology as $BC$ is included.
\end{theorem}
\begin{proof}
Use the mixed-nerve homotopy-pushout description
\[
 BC\simeq\operatorname{hocolim}
 \bigl(BM\longleftarrow E\longrightarrow BN\bigr),
 \qquad
 E=B\bigl(B(S\mathbin{/\mkern-6mu/}N),M,*\bigr).
\]
Lemma~\ref{C:opposite-corner-action-groupoid} is an actual
$M$-equivariant weak equivalence. Applied to the one-point left
$N$-set, the same lemma gives $BN\simeq BG$ through $Bq_N$.
Thus this is equivalently the span with edge
$B(B(eS\mathbin{/\mkern-6mu/}G),M,*)$ and other vertex $BG$.

The groupoid $eS\mathbin{/\mkern-6mu/}G$ has one component for each
$t\in T$. Each component is equivalent to $BG_t$ for a finite
stabilizer subgroup $G_t\leq G$. The map to the component set is
a rational homology equivalence and an $\mathbb F_p$ homology
equivalence when $p\nmid|G|$. Applying the first-quadrant bar
homology spectral sequence to the remaining $M$-bar preserves
these equivalences. The same assertion holds for $BG\to *$.
The map of the two homotopy-pushout spans consequently proves (1);
because the map is the actual functor $Q$, it is an isomorphism
of rings, not merely of graded vector spaces.

For (2) use the full mapping-space form of Miller's theorem A
\cite[Theorem A, p.~39]{miller1984}, with its corrected proof
\cite{miller1985}. For every finite group $H$ and finite CW
complex $X$ the based mapping space $\operatorname{Map}_*(BH,X)$
is weakly contractible. Evaluation and its constant-map section
therefore give
\[
 X\xrightarrow{\simeq}\operatorname{Map}(BH,X).
\]
This is the full mapping-space assertion, not just the
fundamental-group consequence used elsewhere.

The canonical constant-map inclusion
\[
 \operatorname{Map}(T,X)\longrightarrow
 \operatorname{Map}\bigl(B(eS\mathbin{/\mkern-6mu/}G),X\bigr)
\]
is therefore a weak equivalence. It is natural for the actual
right $M$-action; no equivariant choice of component basepoints or
of contractions is needed. Mapping out of a homotopy colimit is
a homotopy limit of these mapping spaces. Equivalently, use the
simplicial $M$-bar and its derived totalization. Levelwise there
are only copies of the preceding weak equivalence, so its
homotopy-limit comparison is a weak equivalence. The corresponding
constant-map inclusion for $BG\to *$ is also a weak equivalence.
Mapping the two pushout spans into $X$ and taking their homotopy
pullbacks proves (2).

Choose a finite CW model $X$ of $BC$. Surjectivity on components
of the mapping-space equivalence lifts the chosen equivalence
$BC\to X$ through $BQ$, up to homotopy. Composing that lift with
an inverse equivalence $X\to BC$ gives (3).
\end{proof}


\begin{corollary}[The precise projective-plane consequences]
\label{C:finite-target-corner-CP2}
In the setting of the preceding theorem, suppose $BC\simeq\mathbb CP^2$.
Write $Y=BC(M,T)$ and let $u$ be an integral generator of
$H^2(BC;\mathbb Z)$. Then $Y$ is simply connected and has the
rational and good-prime cohomology rings of $\mathbb CP^2$.
There is a class $v\in H^2(Y;\mathbb Z)$ with
\[
 (BQ)^*v=u,\qquad (BQ)^*(v^2)=u^2.
\]
The group $H^2(Y;\mathbb Z)$ is infinite cyclic, $v$ is a generator,
and $v^2$ is primitive in the free quotient of $H^4(Y;\mathbb Z)$.
For every finite group $H$ and every map $a:BH\to Y$,
the composite $h a:BH\to BC$ is null-homotopic, and
$a^*v=0$ integrally.
\end{corollary}
\begin{proof}
The quotient functor is surjective on every hom set and fixes the
two objects, so it is surjective on fundamental groups. Since
$BC$ is simply connected, so is $Y$.
Part (1) of the theorem gives its asserted rational and
good-prime cohomology.

Take $v=h^*u$. The homotopy equality $hBQ\simeq1$ proves the
displayed identities integrally. A finite category has finitely
generated homology in each degree. Since $H_1(Y;\mathbb Z)=0$,
the universal coefficient theorem identifies $H^2(Y;\mathbb Z)$
with the dual of $H_2(Y;\mathbb Z)$; its rational rank is one.
It is therefore infinite cyclic. The split surjection $(BQ)^*$
makes $v$ its generator. The same split surjection in degree four,
whose rational rank is one, makes $v^2$ primitive modulo torsion.
Finally Miller's full theorem applies to $h a$ with finite target
$BC\simeq\mathbb CP^2$. Its nullity gives
$a^*v=(h a)^*u=0$.
\end{proof}

\begin{corollary}[A trivial minimal group gives an actual equivalence]
\label{C:trivial-group-opposite-corner}
In Theorem~\ref{C:finite-target-opposite-corner}, suppose the minimal
local group of $N$ is trivial. Then the actual quotient
$BC\to BC(M,N\backslash S)$ is a homotopy equivalence.
Consequently, if $BC\simeq\mathbb CP^2$, then
\[
 |C|\ \geq\ 21+|N|+|S|-|N\backslash S|.
\]
In particular a nontrivial such opposite corner cannot occur in a
projective-plane category with at most $22$ arrows.
\end{corollary}
\begin{proof}
The groupoid in Lemma~\ref{C:opposite-corner-action-groupoid} is
now the discrete set $eS=N\backslash S$, and $BN$ is contractible.
Thus the actual map of mixed-nerve pushout spans in the proof of
Theorem~\ref{C:finite-target-opposite-corner} is an objectwise
homotopy equivalence. This proves the first assertion without a
finite-target or coefficient comparison.
The exact act-cone minimum in
Theorem~\ref{C:act-cone-minimum-twenty-two} gives
$|M|+|N\backslash S|+1\geq22$. Subtracting this size from
$|C|=|M|+|N|+|S|$ gives the stated inequality.
\end{proof}

\begin{lemma}[Finite-target detection after an orbit quotient]
\label{C:finite-target-orbit-wedge}
Let $X=B\mathcal C(M,S)$ satisfy
Lemma~\ref{C:two-column-projective-orbit-kernel}, and suppose $X$
is simply connected. Suppose a map $h:X\to F$ to a simply
connected finite CW complex induces an isomorphism on integral
second cohomology. Write $Y=B\mathcal C(M,\bar S)$ and let $o$
be the number of core orbits. Then there is a map
\[
 Y\longrightarrow F\vee\bigvee_{o-1}S^2
\]
inducing an isomorphism on integral second cohomology.
In particular every integral second-cohomology class on $Y$
restricts to zero along every map $BH\to Y$ from a finite group.
\end{lemma}
\begin{proof}
Corollary~\ref{C:two-column-orbit-wedge} supplies a homotopy
equivalence $X\vee\bigvee_{o-1}S^2\to Y$ extending the actual
quotient on $X$. Compose its homotopy inverse with
$h\vee1:X\vee\bigvee_{o-1}S^2\to F\vee\bigvee_{o-1}S^2$.
The wedge decomposition and the hypothesis on $h^*$ prove the
cohomology assertion. The target is a finite CW complex, so
Miller's theorem makes every composite from $BH$ null-homotopic.
Since every second-cohomology class on $Y$ pulls back from this
target, all the stated restrictions vanish.
\end{proof}

\begin{remark}[Scope of the smaller act cone]
\label{C:finite-target-corner-scope}
Corollary~\ref{C:finite-target-corner-CP2} gives a finite-source
nullity condition for
every local-group restriction of the distinguished class $v$.
It does not identify the integral homotopy type of $Y$.
In particular a lower bound that assumes $Y$ itself has finite
CW type cannot be applied merely from this comparison.
Universal rational cup primitives and certified field cup
obstructions do apply through (1); any imported finite-type hull
argument must separately check that its use of finiteness is
exactly the local-group nullity supplied by the corollary.
\end{remark}



